% !TEX encoding = UTF-8
% !TEX TS-program = pdflatex
% !TEX root = ../tesi.tex

%**************************************************************
\chapter{Introduzione}
\label{cap:introduzione}
%**************************************************************

\intro{Nel seguente capitolo si introduce brevemente l'azienda ospitante e il progetto affrontato.}\\

%Introduzione al contesto applicativo.\\

%\noindent Esempio di utilizzo di un termine nel glossario \\
%\gls{api}. \\

%\noindent Esempio di citazione in linea \\
%\cite{site:agile-manifesto}. \\

%\noindent Esempio di citazione nel pie' di pagina \\
%citazione\footcite{womak:lean-thinking} \\

%**************************************************************
\section{L'azienda}

\noindent{\myCompany}\footnote{Sito ufficiale: \url{https://www.ergon.it/}}(da qui in poi "\textit{Ergon}") è un'azienda italiana, fondata nel 1988, 
situata a Castelfranco Veneto.\\
Essa si occupa principalmente soluzioni gestionali per piccole e medie imprese e dello sviluppo di \textit{software \gls{erpg}} per i settori dell'alimentare e dei trasporti,
ma completa l'offerta con la vendita di prodotti hardware, servizi \textit{web} e \textit{hosting}, nonché con progetti di \textit{\gls{serverconsolidationg}} e virtualizzazione di
sistemi.
L'azienda inoltre si è sviluppata in maniera costante negli anni e oggi può vantare una posizione di tutto rispetto tra le aziende dello stesso settore.
Attualmente fanno parte della stessa gestione:

\begin{itemize}
    \item \textit{Ergon Informatica S.R.L.}: che si occupa del \textit{software};
    \item \textit{Ergon S.R.L.}: che si occupa dei servizi tecnologici;
    \item \textit{Ergon Servizi S.R.L.}: che si occupa dei servizi amministrativi, logistici e di \textit{marketing} delle altre due parti.
\end{itemize}
Il logo dell'azienda è illustrato in \hyperref[fig:logo]{Figura 1.1}.
\begin{figure}[!h]
    \centering
    \includegraphics[width=0.8\columnwidth]{Ergon_Logo.pdf}
    \caption{Logo \myCompany}
    \label{fig:logo}
\end{figure}
\newline
\noindent Il prodotto proprietario dell'azienda è \textit{ERGDIS}, sistema \gls{erpg}
il cui insieme dei moduli copre ogni aspetto della conduzione aziendale.\newline

\noindent In particolare vengono gestiti vari compiti che si dislocano dall'area amministrativa al controllo direzionale,
dall'area commerciale alla pianificazione e al controllo della produzione, dalla gestione acquisti alla logistica di magazzino, 
all'archiviazione ottica alla gestione della qualità.\\
Alcuni di essi, inoltre, si possono interfacciare con dispositivi automatici presenti in azienda, come, ad esempio,
linee di confezionamento o \textit{robot}.\\

Paragrafo riguardo alle funzionalità di ERGDIS\\

Paragrafo riguardo agli obiettivi generali dell'azienda
 

%**************************************************************
\section{L'idea}
\noindent Lo \textit{stage} proposto consiste nella progettazione e nello sviluppo di un modulo \textit{software} volto
ad assistere l'azienda nella fase di approvvigionamento dei prodotti dai propri fornitori, supportandola nel scegliere
da quale fornitore e quando acquistare i prodotti.\\

\noindent Questa nuova funzionalità andrebbe ad ottimizzare un modulo già esistente che però, per ogni prodotto da ordinare, prende in
considerazione l'ultima data d'ordine disponibile prima dell'inizio
dell'effettiva copertura del fabbisogno del prodotto stesso. Questo dunque non garantirebbe con certezza una scelta ottimale
in relazione alle possibilità d'ordine fornite dagli appositi listini e calendario dei fornitori.\\
Data la natura combinatoria del problema, il modulo dovrà fornire in tempi ragionevoli una "buona soluzione"
del problema, ovvero tendente il più possibile all'ottimo, e dovrà integrarsi con l'intero sistema \textit{ERGDIS}.\\

\noindent È previsto inoltre che i dati su cui si è fatta l'ottimizzazione e il confronto dei risultati vengano visualizzati tramite un'apposita
interfaccia grafica che verrà sviluppata in linea con l'ambiente di sviluppo dell'azienda (\textit{\gls{.netg}} e \textit{\gls{devexpressg}}).


%**************************************************************
\section{Descrizione dello stage}

%**************************************************************
\subsection{Introduzione}

%**************************************************************
\subsection{Obiettivi}

\noindent Di seguito vengono elencati tutti gli obiettivi previsti dallo \textit{stage}:
\begin{itemize}
    \item Analisi del contesto \textit{ERGDIS};
    \item Studio dei principali algoritmi di ricerca operativa e ottimizzazione combinatoria;
    \item Redazione di uno studio di fattibilità con integrazione di micro-moduli di test;
    \item Redazione di un'analisi dei requisiti;
    \item Sviluppo e codifica del modulo \textit{software} con le tecnologie utilizzate dall'azienda;
    \item Redazione di documentazione tecnica riguardante le scelte implementative e architetturali effettuate;
    \item \textit{Report} finale sui risultati ottenuti.
\end{itemize}
%**************************************************************
\subsection{Analisi preventiva dei rischi}
Durante la fase iniziale dello \textit{stage}, sono stati rilevati dei possibili rischi che avrebbero potuto presentarsi
durante il percorso del progetto.\\
Si sono dunque trovate delle soluzioni che potessero arginare i problemi. In particolare:
\begin{enumerate}
    \item \textbf{Comprensione e confronto degli algoritmi}\\[0.2cm]
    \textbf{Problema:} il progetto richiede un'ampia fase di studio che riguarda principalmente la teoria
    di tecniche per la risoluzione di problemi di ricerca operativa e ottimizzazione combinatoria.
    Questo poteva portare a presentarsi la possibilità di non comprendere fino in fondo l'algoritmo, e poteva essere difficile cogliere e confrontare
    i pregi e difetti di ciascuno di essi.\\[0.2cm]
    \textbf{Soluzione:} è stato organizzato un incontro iniziale con il tutor per fornire una base da cui poi iniziare una ricerca più approfondita.
    Sono stati forniti anche delle dispense utili per rafforzare la base di partenza.
    \item \textbf{Tecnologie e ambiente di sviluppo}\\[0.2cm]
    \textbf{Problema:} venivano richieste alcune tecnologie, come per esempio \textit{Entity Framework} o \textit{\gls{devexpressg}} a me assolutamente ignote.
    Sebbene avessi delle basi abbastanza solide di \textit{C\#} derivanti dalla conoscenza di altri linguaggi quali \textit{C++} e \textit{Java}, venivano richieste tuttavia
    alcune tecnologie integrate nel linguaggio (\textit{LINQ}) anch'esse ignote. L'ambiente di sviluppo e l'\textit{IDE} non erano mai stati utilizzati.\\[0.2cm]
    \textbf{Soluzione:} sono stati forniti dei riferimenti consigliati per l'autoapprendimento. Tuttavia qualsiasi dubbio ragionevolmente particolare
    poteva essere richiesto al tutor. È stato effettuato insieme al tutor il \textit{setup} dell'ambiente di sviluppo e la conseguente creazione dei \textit{database}.
    \item \textbf{Calibrazione dei parametri e funzione obiettivo}\\[0.2cm]
    \textbf{Problema:} dopo la scelta e l'implementazione dell'algoritmo, è molto importante:
    \begin{itemize}
        \item definire una funzione obiettivo che vada a descrivere in maniera "buona" l'andamento dell'algoritmo stesso;
        \item calibrare i parametri in base allo spazio delle soluzioni del problema preso in esame.
    \end{itemize}
    Tuttavia entrambe sono azioni molto delicate che possono compromettere il funzionamento stesso dell'algoritmo anche
    se implementato correttamente.\\[0.2cm]
    \textbf{Soluzione:} cercare una costruzione e calibrazione per passi ed empiricamente dimostrabili e presentarle in una discussione con il tutor,
    in modo tale da creare una \textit{baseline} su cui basarsi per continuare con i passi successivi.


\end{enumerate}

\noindent
%**************************************************************
\section{Organizzazione del testo}
\noindent Di seguito viene illustrata l'organizzazione dei capitoli successivi:
\begin{description}
    \item[{\hyperref[cap:studio-fattibilita]{Il secondo capitolo}}] 
    approfondice lo studio di fattibilità effettuato, utile per entrare a conoscenza
    delle più utilizzate tecniche di ottimizzazione combinatoria e per analizzare
    quali siano i vantaggi e svantaggi di ognuno di essi.
    
    \item[{\hyperref[cap:analisi-requisiti]{Il terzo capitolo}}]
    descrive l'analisi dei requisiti del progetto, comprensiva di diagrammi dei
    casi d'uso e raccolta dei requisiti derivanti dall'analisi di questi ultimi.
    
    \item[{\hyperref[cap:progettazione-codifica]{Il quarto capitolo}}]
    approfondisce le fasi di progettazione e codifica, comprensiva di diagrammi delle classi
    e di approfondimenti a livello implementativo.
    
    \item[{\hyperref[cap:verifica-validazione]{Il quinto capitolo}}]
    espone tutte verifiche effettuate durante il progetto e la validazione
    finale a conferma dei requisiti inizialmente stilati nella fase di
    analisi dei requisiti.
    
    \item[{\hyperref[cap:conclusioni]{Il sesto capitolo}}]
    presenta le conclusione tratte dallo \textit{stage}, comprensivo di conoscenze
    acquisite e considerazioni di carattere personale.
\end{description}

\noindent Riguardo la stesura del testo, relativamente al documento sono state adottate le seguenti convenzioni tipografiche:
\begin{itemize}
	\item gli acronimi, le abbreviazioni e i termini ambigui o di uso non comune menzionati vengono definiti nel glossario, situato alla fine del presente documento;
	\item i termini in lingua straniera o facenti parti del gergo tecnico sono evidenziati con il carattere \emph{corsivo}.
\end{itemize}